\chapter{Diseño del controlador}\label{chap:controlador}

El capítulo~\ref{chap:modelo} dejó establecida la planta y sus límites: un equilibrio inestable, un actuador cuyo voltaje alcanza como máximo $3{,}5$~V y un intervalo de posiciones en el que el imán puede moverse de verdad. Este capítulo responde a la pregunta siguiente, que es qué controlador lineal estabiliza esa planta y con qué holgura. Se linealiza la planta en el punto de diseño y se fijan los objetivos, se presenta el controlador con doble acción integral (PII), se examina su robustez en el intervalo de operación y, por último, se define el controlador de comparación con una sola acción integral. Cuánto logra el PII frente a las no linealidades no suaves se evalúa en los capítulos~\ref{chap:desempeno} y~\ref{chap:doble_integral}.

\section{Planta linealizada y objetivos del diseño}\label{sec:planta_diseno}

El diseño parte del modelo linealizado~\eqref{eq:lineal} del capítulo~\ref{chap:modelo}, evaluado en el punto de equilibrio $x_1^*=-4$~cm. La entrada es el voltaje del actuador y la salida es la posición del imán. Con esos datos, la función de transferencia de la planta es
\begin{equation}
G(s)=\frac{281{,}7}{s^{2}+8{,}563\,s-351{,}2} .
\label{eq:planta}
\end{equation}
El numerador $281{,}7$ es la ganancia de entrada $1/\bigl(mb(a-x_1^*)^4\bigr)$ del modelo~\eqref{eq:lineal}, y el término independiente del denominador, $-351{,}2$, es el valor de $4g/(a-x_1^*)$ de la ecuación característica~\eqref{eq:caracteristica}. Su signo negativo produce un polo real en el semiplano derecho, en $14{,}94$~rad/s, que es la inestabilidad en lazo abierto de la configuración atractiva descrita en el capítulo~\ref{chap:introduccion}. El coeficiente $8{,}563$ es la fricción viscosa $c_1$.

El controlador debe cumplir cuatro objetivos. Debe estabilizar el lazo, para lo cual el diagrama de Nyquist ha de rodear el punto $-1$ una vez en sentido antihorario, como exige el criterio con un polo inestable (capítulo~\ref{chap:fundamentos}). Debe dejar márgenes de estabilidad razonables. Debe anular el error estacionario ante escalones y rampas, lo que pide un lazo de tipo dos (capítulo~\ref{chap:fundamentos}). Y debe proporcionar alta ganancia a baja frecuencia, la motivación de la doble acción integral expuesta en la sección~\ref{sec:motivacion}.

\section{El controlador con doble acción integral}\label{sec:pii}

El controlador se ajustó sobre la planta~\eqref{eq:planta}, de forma interactiva en los editores de Bode y del lugar de las raíces del \emph{Control System Designer} de MATLAB. El resultado es
\begin{equation}
C(s)=\frac{152{,}43\,(s+4{,}875)(s+1{,}861)(s+17{,}31)}
{s^{2}(s+301{,}1)} .
\label{eq:pii}
\end{equation}

\subsection{Estructura}

El controlador~\eqref{eq:pii} tiene tres partes. El doble integrador $1/s^2$ hace al lazo de tipo dos, lo que anula el error estacionario ante escalones y rampas, y proporciona la acción acumulada que se espera que ayude a salir de la zona muerta. Los ceros en $-1{,}861$ y $-4{,}875$~rad/s compensan el atraso de fase de los dos integradores por debajo de la frecuencia de cruce. El par cero-polo en $-17{,}31$ y $-301{,}1$~rad/s actúa como red de adelanto, con su máximo aporte de fase cerca de $72$~rad/s.

\subsection{Márgenes de estabilidad}

Con esta estructura, el lazo abierto $L(s)=C(s)G(s)$ tiene un margen de fase de $59{,}9^\circ$ en $129{,}3$~rad/s (figura~\ref{bode_fase}). Como la planta tiene un polo inestable, el margen de ganancia es inferior: el lazo sigue siendo estable si la ganancia aumenta, pero se inestabiliza si disminuye por debajo de $0{,}14$ veces su valor nominal ($-17{,}1$~dB en $9{,}63$~rad/s). Los polos del lazo cerrado son $-143{,}3\pm133{,}7\mathrm{j}$, $-14{,}10$, $-7{,}28$ y $-1{,}71$, todos en el semiplano izquierdo.

\begin{figure}[ht]
    \centering
    % Estilo común de las figuras de diseño del PII (pgfplots).
% Generado por diseno_PII_tikz.m
\def\PIIwg{9.6320}
\def\PIImgwg{17.0713}
\def\PIIwp{129.2942}
\def\PIIphwp{-120.1396}
\def\PIIgmdb{-17.1}
\def\PIIpm{59.9}
\def\PIIos{16.9}
\def\PIIts{0.284}
\def\PIItp{0.0238}
\def\PIIyp{1.1688}
\def\PIIflechaA{(-7.6265,1.7150) -- (-7.5765,1.4566)}
\def\PIIflechaC{(-7.5765,-1.4566) -- (-7.6265,-1.7150)}
\def\PIIflechaB{(-2.6114,-3.2049) -- (-2.5176,-3.1809)}
\def\PIIflechaD{(-2.5176,3.1809) -- (-2.6114,3.2049)}

\pgfplotsset{
  diseno/.style={
    grid=major, grid style={gray!25},
    tick label style={font=\small}, label style={font=\small}, legend style={font=\small},
    /pgf/number format/use comma,
    table/col sep=comma,
    every axis plot/.append style={line join=round},
  },
}
\definecolor{tray}{RGB}{31,94,166}
\definecolor{marg}{RGB}{196,96,40}
\definecolor{refc}{RGB}{40,40,40}

\begin{tikzpicture}
\begin{axis}[diseno, name=mag, width=0.85\textwidth, height=5.2cm, xmode=log,
             xmin=1e-2, xmax=1e4, ymin=-100, ymax=110, ytick={-100,-50,0,50,100},
             xticklabels={}, ylabel={magnitud [\si{dB}]}]
  \addplot[refc!50, densely dashed] coordinates {(1e-2,0) (1e4,0)};
  \addplot[tray, line width=1pt] table[x=w, y=mag_db] {images/pii_design/tikz/bode.csv};
  \draw[marg, line width=0.8pt] (axis cs:\PIIwg,0) -- (axis cs:\PIIwg,\PIImgwg);
  \addplot[marg, only marks, mark=*, mark size=1.8pt] coordinates {(\PIIwg,\PIImgwg)};
  \node[font=\small, anchor=north east, align=left, fill=white, inner sep=2pt]
    at (axis cs:8e3,104) {MG $=\num{\PIIgmdb}$~\si{dB} en $\omega=\num[round-mode=places,round-precision=2]{\PIIwg}$~\si{rad/s}};
\end{axis}
\begin{axis}[diseno, at={(mag.south west)}, anchor=north west, yshift=-0.35cm,
             width=0.85\textwidth, height=5.2cm, xmode=log,
             xmin=1e-2, xmax=1e4, ymin=-375, ymax=-90, ytick={-360,-315,-270,-225,-180,-135},
             xlabel={frecuencia [\si{rad/s}]}, ylabel={fase [\si{\degree}]}]
  \addplot[refc!50, densely dashed] coordinates {(1e-2,-180) (1e4,-180)};
  \addplot[tray, line width=1pt] table[x=w, y=fase] {images/pii_design/tikz/bode.csv};
  \draw[marg, line width=0.8pt] (axis cs:\PIIwp,-180) -- (axis cs:\PIIwp,\PIIphwp);
  \addplot[marg, only marks, mark=*, mark size=1.8pt] coordinates {(\PIIwp,\PIIphwp)};
  \node[font=\small, anchor=south east, align=left, fill=white, inner sep=2pt]
    at (axis cs:8e3,-370) {MF $=\num{\PIIpm}\si{\degree}$ en $\omega=\num[round-mode=places,round-precision=1]{\PIIwp}$~\si{rad/s}};
\end{axis}
\end{tikzpicture}

    \caption[Diagrama de Bode]{Diagrama de Bode del lazo abierto $L(s)=C(s)G(s)$ en $x_1^*=-4$~cm. Se indican el margen de ganancia (inferior, por el polo inestable de la planta) y el margen de fase.}
    \label{bode_fase}
\end{figure}

El diagrama de Nyquist (figura~\ref{nyquist}) rodea una vez el punto $-1$ en sentido antihorario, que es el rodeo que exige el criterio con un polo inestable en lazo abierto.

\begin{figure}[ht]
    \centering
    % Estilo común de las figuras de diseño del PII (pgfplots).
% Generado por diseno_PII_tikz.m
\def\PIIwg{9.6320}
\def\PIImgwg{17.0713}
\def\PIIwp{129.2942}
\def\PIIphwp{-120.1396}
\def\PIIgmdb{-17.1}
\def\PIIpm{59.9}
\def\PIIos{16.9}
\def\PIIts{0.284}
\def\PIItp{0.0238}
\def\PIIyp{1.1688}
\def\PIIflechaA{(-7.6265,1.7150) -- (-7.5765,1.4566)}
\def\PIIflechaC{(-7.5765,-1.4566) -- (-7.6265,-1.7150)}
\def\PIIflechaB{(-2.6114,-3.2049) -- (-2.5176,-3.1809)}
\def\PIIflechaD{(-2.5176,3.1809) -- (-2.6114,3.2049)}

\pgfplotsset{
  diseno/.style={
    grid=major, grid style={gray!25},
    tick label style={font=\small}, label style={font=\small}, legend style={font=\small},
    /pgf/number format/use comma,
    table/col sep=comma,
    every axis plot/.append style={line join=round},
  },
}
\definecolor{tray}{RGB}{31,94,166}
\definecolor{marg}{RGB}{196,96,40}
\definecolor{refc}{RGB}{40,40,40}

\begin{tikzpicture}
\begin{axis}[diseno, width=0.8\textwidth, height=0.62\textwidth,
             xmin=-10, xmax=2, ymin=-7, ymax=7, axis equal image,
             xlabel={$\operatorname{Re}\{L(j\omega)\}$}, ylabel={$\operatorname{Im}\{L(j\omega)\}$},
             legend pos=north west, legend cell align=left]
  \addplot[refc!40, thin, forget plot] coordinates {(-10,0) (2,0)};
  \addplot[refc!40, thin, forget plot] coordinates {(0,-7) (0,7)};
  \addplot[refc!55, densely dashed, domain=0:360, samples=120, forget plot] ({cos(x)},{sin(x)});
  \addplot[tray, line width=1pt]
    table[x=re, y=im] {images/pii_design/tikz/nyquist_pos.csv};
  \addlegendentry{$\omega>0$}
  \addplot[tray!55, line width=0.8pt, densely dashed]
    table[x=re, y=im] {images/pii_design/tikz/nyquist_neg.csv};
  \addlegendentry{$\omega<0$}
  \addplot[red!75!black, only marks, mark=+, mark size=4pt, line width=1.2pt] coordinates {(-1,0)};
  \addlegendentry{$-1$}
  \draw[tray, -stealth, line width=1pt] \PIIflechaA;
  \draw[tray, -stealth, line width=1pt] \PIIflechaB;
  \draw[tray!55, -stealth, line width=0.8pt] \PIIflechaC;
  \draw[tray!55, -stealth, line width=0.8pt] \PIIflechaD;
  \draw[marg, thin] (axis cs:-7.14,0) -- (axis cs:-8.5,-1.2);
  \node[font=\small, anchor=north, fill=white, inner sep=1.5pt] at (axis cs:-8.5,-1.2) {$-1/\mathrm{MG}$};
  \addplot[marg, only marks, mark=*, mark size=1.6pt, forget plot] coordinates {(-7.14,0)};
\end{axis}
\end{tikzpicture}

    \caption[Diagrama de Nyquist]{Diagrama de Nyquist de $L(s)$. Con un polo de lazo abierto en el semiplano derecho ($P=1$), la estabilidad en lazo cerrado exige un rodeo antihorario al punto $-1$.}
    \label{nyquist}
\end{figure}

\subsection{Respuesta al escalón}

En lazo cerrado lineal, la respuesta al escalón (figura~\ref{step_response}) tiene un tiempo de subida de $9{,}3$~ms, un sobrepaso de $16{,}9$\,\% con pico de $1{,}169$ en $23{,}8$~ms y un tiempo de establecimiento de $0{,}284$~s.

\begin{figure}[ht]
    \centering
    % Estilo común de las figuras de diseño del PII (pgfplots).
% Generado por diseno_PII_tikz.m
\def\PIIwg{9.6320}
\def\PIImgwg{17.0713}
\def\PIIwp{129.2942}
\def\PIIphwp{-120.1396}
\def\PIIgmdb{-17.1}
\def\PIIpm{59.9}
\def\PIIos{16.9}
\def\PIIts{0.284}
\def\PIItp{0.0238}
\def\PIIyp{1.1688}
\def\PIIflechaA{(-7.6265,1.7150) -- (-7.5765,1.4566)}
\def\PIIflechaC{(-7.5765,-1.4566) -- (-7.6265,-1.7150)}
\def\PIIflechaB{(-2.6114,-3.2049) -- (-2.5176,-3.1809)}
\def\PIIflechaD{(-2.5176,3.1809) -- (-2.6114,3.2049)}

\pgfplotsset{
  diseno/.style={
    grid=major, grid style={gray!25},
    tick label style={font=\small}, label style={font=\small}, legend style={font=\small},
    /pgf/number format/use comma,
    table/col sep=comma,
    every axis plot/.append style={line join=round},
  },
}
\definecolor{tray}{RGB}{31,94,166}
\definecolor{marg}{RGB}{196,96,40}
\definecolor{refc}{RGB}{40,40,40}

\begin{tikzpicture}
\begin{axis}[diseno, width=0.85\textwidth, height=6cm, xmin=0, xmax=0.6, ymin=0, ymax=1.32,
             xlabel={tiempo [\si{s}]}, ylabel={amplitud},
             xtick={0,0.1,...,0.6}, x tick label style={/pgf/number format/fixed}]
  \addplot[refc!55, densely dashed] coordinates {(0,1) (0.6,1)};
  \addplot[refc!30, densely dotted, forget plot] coordinates {(0,0.98) (0.6,0.98)};
  \addplot[refc!30, densely dotted, forget plot] coordinates {(0,1.02) (0.6,1.02)};
  \addplot[tray, line width=1pt] table[x=t, y=y] {images/pii_design/tikz/escalon.csv};
  \addplot[marg, only marks, mark=*, mark size=1.8pt] coordinates {(\PIItp,\PIIyp)};
  \node[font=\small, anchor=south west, inner sep=2pt] at (axis cs:\PIItp,\PIIyp)
    {sobrepaso \SI{\PIIos}{\percent}};
  \draw[marg, densely dashed] (axis cs:\PIIts,0) -- (axis cs:\PIIts,0.98);
  \node[font=\small, anchor=south west, fill=white, inner sep=2pt] at (axis cs:\PIIts,0.08)
    {$t_s=\num{\PIIts}$~\si{s} (\SI{2}{\percent})};
\end{axis}
\end{tikzpicture}

    \caption[Respuesta al escalón]{Respuesta al escalón unitario del lazo cerrado linealizado.}
    \label{step_response}
\end{figure}

\section{Robustez en el intervalo de operación}\label{sec:robustez_pii}

El controlador se diseñó en $-4$~cm, pero las pruebas se hacen entre $-3$ y $-2$~cm. La sección~\ref{sec:actuador_estatico} del capítulo~\ref{chap:modelo} justifica ese intervalo por el voltaje que el actuador puede entregar, y la sección~\ref{sec:intervalo} del capítulo~\ref{chap:desempeno} lo confirma con un barrido de escalones. Conviene entonces saber si el diseño se sostiene donde se opera. Al acercarse el imán a la bobina, el polo inestable se desplaza hacia la derecha, de $14{,}94$ a $16{,}84$~rad/s, y la planta se vuelve más difícil de estabilizar. La tabla~\ref{tab:pii_intervalo} reúne el margen de fase y el sobrepaso lineal del mismo controlador~\eqref{eq:pii}, sin reajuste, sobre la planta linealizada en cada posición.

\begin{table}[ht]
    \centering
    \caption[PII en el intervalo de operación]{Controlador PII~\eqref{eq:pii}, diseñado en $-4$~cm, evaluado sobre la planta linealizada en las posiciones de prueba.}
    \label{tab:pii_intervalo}
    \begin{tabular}{rccc}
        \hline
        $x_1^*$ [cm] & Polo inestable [rad/s] & Margen de fase & Sobrepaso lineal \\
        \hline
        $-4{,}0$ (diseño) & $14{,}94$ & $59{,}9^\circ$ & $16{,}9$\,\% \\
        $-3{,}0$ & $15{,}82$ & $54{,}5^\circ$ & $19{,}2$\,\% \\
        $-2{,}5$ & $16{,}31$ & $51{,}1^\circ$ & $21{,}6$\,\% \\
        $-2{,}0$ & $16{,}84$ & $47{,}4^\circ$ & $24{,}7$\,\% \\
        \hline
    \end{tabular}
\end{table}

El lazo permanece estable en todo el intervalo, con menos margen conforme el imán se acerca a la bobina. El margen de fase baja de $59{,}9^\circ$ a $47{,}4^\circ$ y el sobrepaso lineal sube de $16{,}9$\,\% a $24{,}7$\,\%. El controlador, por tanto, no necesita reajustarse para operar entre $-3$ y $-2$~cm, y todas las pruebas de los capítulos siguientes lo emplean tal como está.

\section{El controlador PI de comparación}\label{sec:pi}

La pregunta central del trabajo es qué aporta la segunda integración. Para responderla se necesita un controlador de referencia que difiera del PII en eso y en poco más. El controlador PI de comparación es
\begin{equation}
C_{PI}(s)=\frac{178{,}14\,(s+15)(s+12{,}91)}{s\,(s+400)} .
\label{eq:pi}
\end{equation}
A diferencia del PII~\eqref{eq:pii}, tiene un solo integrador, de modo que el lazo es de tipo uno y su ganancia a baja frecuencia crece más lentamente. Conserva la red de adelanto, con un polo en $-400$~rad/s, y dos ceros en $-15$ y $-12{,}91$~rad/s. Sobre la planta~\eqref{eq:planta} tiene un margen de fase de $64{,}1^\circ$ en $118{,}6$~rad/s, un sobrepaso lineal de $17{,}6$\,\% y un tiempo de establecimiento lineal de $0{,}240$~s. Estos valores son comparables a los del PII ($59{,}9^\circ$, $16{,}9$\,\% y $0{,}284$~s), de modo que la diferencia entre ambos lazos queda en la segunda acción integral y no en una respuesta lineal mejor o peor sintonizada.

Este controlador servirá en el capítulo~\ref{chap:doble_integral} para aislar el efecto de la segunda integración sobre el comportamiento con atascamiento-deslizamiento y zona muerta. Antes, el capítulo~\ref{chap:desempeno} evalúa el desempeño del PII en simulación.
